% ==============================================================================
% dxh_toolset.tex — DXH workshop: today's toolset (teacher-friendly ecosystem)
% DXH講座: 今日つかう道具一式（高校教員向けに平易化したエコシステム図）
%
% The repo-directory ecosystem_layers diagram is developer-oriented and
% meaningless to high-school teachers; this shows the same idea as the three
% physical tools on the desk plus the freely published materials.
% リポジトリ構成ベースの ecosystem_layers は開発者向けで高校教員には伝わらない
% ため、机上の3つの道具＋無料公開教材という形で同じ趣旨を示す。
% ==============================================================================
\documentclass[border=10pt]{standalone}
\usepackage{tikz}
\usetikzlibrary{arrows.meta, positioning}
\usepackage{luatexja}
\usepackage[match]{luatexja-fontspec}
% Cross-platform font: Hiragino (macOS) → Yu Gothic (Windows) → Noto (Linux)
\usepackage{ifthen}
\newboolean{jfontset}\setboolean{jfontset}{false}
\IfFontExistsTF{Hiragino Kaku Gothic ProN}{%
  \setmainjfont{Hiragino Kaku Gothic ProN}%
  \setsansjfont{Hiragino Kaku Gothic ProN}%
  \setboolean{jfontset}{true}}{}
\ifthenelse{\boolean{jfontset}}{}{%
  \IfFontExistsTF{Yu Gothic}{%
    \setmainjfont{Yu Gothic}%
    \setsansjfont{Yu Gothic}%
    \setboolean{jfontset}{true}}{}}
\ifthenelse{\boolean{jfontset}}{}{%
  \IfFontExistsTF{Noto Sans CJK JP}{%
    \setmainjfont{Noto Sans CJK JP}%
    \setsansjfont{Noto Sans CJK JP}}{}}

\begin{document}

% Colors defined after \begin{document} (standalone compatibility, rule T9)
% 色定義は \begin{document} の後（standalone 互換, ルール T9）
\definecolor{cpc}{RGB}{0, 100, 180}
\definecolor{cfly}{RGB}{180, 60, 40}
\definecolor{cjoy}{RGB}{40, 140, 60}
\definecolor{copen}{RGB}{120, 80, 160}

\begin{tikzpicture}[
    tool/.style={
      rectangle, rounded corners=5pt, draw=#1, line width=1.5pt,
      fill=#1!10, minimum width=44mm, minimum height=20mm,
      font=\bfseries\large, text=#1, align=center
    },
    desc/.style={
      font=\small, text=black!60, align=center
    },
    arr/.style={
      -{Stealth[length=3.5mm]}, line width=1.5pt, black!60
    },
    arrlabel/.style={
      font=\small, text=black!55, align=center
    },
  ]

  % === Three tools on the desk / 机の上の3つの道具 ===
  \node[tool=cpc]  (pc)  at (-6.6, 0) {ノートPC\\[-1pt]{\normalsize プログラムを書く}};
  \node[tool=cfly] (fly) at (0, 0)    {StampFly 機体\\[-1pt]{\normalsize プログラムで飛ぶ}};
  \node[tool=cjoy] (joy) at (6.6, 0)  {コントローラ\\[-1pt]{\normalsize スティックで操縦}};

  \node[desc, below=2mm of pc]  {開発環境はインストール済み};
  \node[desc, below=2mm of fly] {センサ＋小型コンピュータ搭載};
  \node[desc, below=2mm of joy] {Atom JoyStick};

  % === Horizontal connections / 横方向の接続（直角ルーティング: 水平のみ） ===
  \draw[arr] (pc.east) -- (fly.west)
    node[arrlabel, midway, above=1mm] {USBで\\書き込み};
  \draw[arr] (joy.west) -- (fly.east)
    node[arrlabel, midway, above=1mm] {電波で\\操縦};

  % === Everything is published for free / すべて無料公開 ===
  \node[rectangle, rounded corners=5pt, draw=copen, line width=1.5pt,
        fill=copen!8, minimum width=152mm, minimum height=15mm,
        align=center, below=13mm of fly, text=copen]
    {\bfseries\large ソースコード・教材・シミュレータはすべて Web で無料公開\\[0pt]
     {\normalsize\mdseries\color{black!60} 学校でも今日と同じ環境をそのまま用意できる}};

\end{tikzpicture}

\end{document}
